\documentclass[11pt]{article}
\usepackage[margin=25mm]{geometry}
\usepackage{amsmath,amssymb}
\begin{document}

\title{I01 prevalidation 9 result:\\Paired coherent reference access produces constraint leakage}
\date{2026-10-04}
\maketitle

\section*{Question and intervention}
The stage-8 one-sided $R_x(\epsilon)$ fault leaked amplitude out of the reference
sector but did not make the $P=0$ pooled constraint non-null. This experiment used a
matched reference intervention: $R_x(\epsilon)$ immediately after preparation and
$R_x(-\epsilon)$ immediately after the scientific phase layer, both on qubit zero.
It asks whether coherent access on both sides of the phase layer permits
recombination into the reference sector.

\section*{Exact Aer result}
For the skewed allocation and $P=0$, $\tau=0$, the maximum pooled joint-probability
difference is $2.04\times10^{-7}$, $1.83\times10^{-6}$, and $2.03\times10^{-5}$ at
$\epsilon=0.01,0.03,0.10$ radians. The approximately quadratic growth occurs while
the unrestricted distributional null remains at numerical zero. Thus the effect is
not a generic label-allocation artefact: it requires matched coherent reference
access and appears in the previously exact constraint.

\section*{Interpretation}
Leaving a sector is insufficient. A later coherent operation must return amplitude
so that the all-zero event can interfere with it. At $\epsilon=0.10$, finite-shot
noncentrality is only $2.14\times10^{-3}$, so this is a mechanism confirmation, not
a finite-shot detection claim.

\section*{Reproduction}
Implementation: \texttt{../src/relational\_atlas\_i01/holdout/qiskit\_rank\_holdout9\_paired\_coherence.py}.\\
Exact output: \texttt{../results/qiskit\_rank\_holdout9\_paired\_coherence.json}.
\end{document}
